\chapter*{Resum}
\addcontentsline{toc}{chapter}{Resum}

Synapse Notes és una plataforma SaaS multi-tenant per a la gestió de notes
personals amb xat contextual basat en Retrieval-Augmented Generation (RAG).
Aquest TFG amplia la plataforma amb un servidor Model Context Protocol
(MCP) que exposa vuit eines de lectura i escriptura de notes amb
autenticació OAuth~2.1, dos agents en segon pla (auto-tag i weekly-digest)
desplegats com a Vercel Cron, i un model d'amenaces explícit basat en el
marc de la Lethal Trifecta. La contribució principal és una anàlisi
empírica de l'aïllament multi-tenant (Row-Level Security), el rendiment
de la cerca vectorial amb pgvector i la resistència a la injecció
indirecta de prompt mesurada amb una suite automatitzada de red-team
(Promptfoo).

\selectlanguage{spanish}
\chapter*{Resumen}
\addcontentsline{toc}{chapter}{Resumen}

Synapse Notes es una plataforma SaaS multi-tenant para la gestión de
notas personales con chat contextual basado en Retrieval-Augmented
Generation (RAG). Este TFG amplía la plataforma con un servidor Model
Context Protocol (MCP) que expone ocho herramientas de lectura y
escritura de notas con autenticación OAuth~2.1, dos agentes en segundo
plano (auto-tag y weekly-digest) desplegados como Vercel Cron, y un
modelo de amenazas explícito basado en el marco de la Lethal Trifecta.
La contribución principal es un análisis empírico del aislamiento
multi-tenant (Row-Level Security), el rendimiento de la búsqueda
vectorial con pgvector y la resistencia a la inyección indirecta de
prompt medida mediante una suite automatizada de red-team (Promptfoo).

\selectlanguage{english}
\chapter*{Abstract}
\addcontentsline{toc}{chapter}{Abstract}

Synapse Notes is a multi-tenant SaaS platform for personal note
management with a Retrieval-Augmented Generation (RAG) chat interface.
This thesis extends the platform with a Model Context Protocol (MCP)
server exposing eight read and write tools over a user's notes with
OAuth~2.1 authentication, two background agents (auto-tag and
weekly-digest) deployed as Vercel Cron jobs, and an explicit threat
model grounded in the Lethal Trifecta framework. The main contribution
is an empirical evaluation of multi-tenant isolation (Row-Level
Security), vector search performance with pgvector, and resistance to
indirect prompt injection measured through an automated red-team suite
(Promptfoo).

\selectlanguage{catalan}
